\documentclass[../../../include/open-logic-section]{subfiles}
%READERNOTE{SOL6-A234: संच-फलनाच्या प्रतिमेत पदमूल्य वापरताना घटक फलनाच्या प्रांतात असावा ही अट स्पष्ट केली. मनमाना संच घेतल्यास प्रतिमा फलनाच्या आलेखावरून परिभाषित करून विभक्तीकरणाने मिळते; यासाठी प्रतिस्थापन लागत नाही.}

\begin{document}

\olfileid{sth}{ordinals}{replacement}
\olsection{प्रतिस्थापन}

\olref[sth][ordinals][vn]{sec} मध्ये क्रमसूचक संख्यांना क्रमप्रकार म्हणून, म्हणजे सुक्रमणांचे प्रमाणभूत प्रतिनिधी म्हणून, घेता येईल असे सुचवून त्यांची ओळख करून दिली. हे साधण्यासाठी \emph{प्रत्येक सुक्रमण कोणत्या तरी क्रमसूचक संख्येशी समरूपी आहे} हे सिद्ध करावे लागेल. मग $\tuple{A, <} \isomorphic \alpha$ अशी क्रमसूचक संख्या $\alpha$ घेऊन $\ordtype{A, <}$ परिभाषित करता येईल.

दुर्दैवाने, आतापर्यंत आपण दिलेल्या स्वयंसिद्धकांवरून हा अपेक्षित निष्कर्ष \emph{सिद्ध करता येत नाही}. का ते \olref[replacement][strength]{sec} मध्ये पाहू; सध्या इतके लक्षात ठेवू की तो सिद्ध होत नाही. म्हणून आपल्याला एक नवे तत्त्व हवे आहे:

\begin{axiom}[प्रतिस्थापन-रूपबंध]
कोणत्याही सूत्र $\phi(x, y)$ साठी पुढील विधान हे एक स्वयंसिद्धक आहे:
\begin{quote}
	कोणत्याही $A$ साठी, $(\forall x \in A)\lexists![y][\phi(x,y)]$ असेल, तर $\Setabs{y}{(\exists x \in A)\phi(x,y)}$ अस्तित्वात असतो.
\end{quote}
\end{axiom}
\noindent
विभक्तीकरणाप्रमाणे हाही एक रूपबंध आहे: अनंत संख्येने असलेल्या निरनिराळ्या $\phi$ साठी त्यातून अनंत संख्येने स्वयंसिद्धके मिळतात. हेच तत्त्व पुढीलप्रमाणेही लिहिता येते, आणि सामान्यतः तसेच लिहिले जाते:

\begin{defish}
``$B$'' नसलेल्या कोणत्याही सूत्र $\phi(x,y)$ साठी पुढील विधान हे एक स्वयंसिद्धक आहे:
\[
\forall A[(\forall x \in A)\lexists![y][\phi(x,y)] \lif \exists B\forall y (y \in B \liff (\exists x \in A)\phi(x,y))]
\]
\end{defish}

पहिल्यांदा पाहताना हे सूत्र गुंतागुंतीचे वाटते. प्रतिस्थापनाचा पुढील सोपा परिणाम मात्र त्यामागील कल्पना \emph{अधिक स्पष्टपणे} व्यक्त करतो:\footnote{पद म्हणजे अशी अभिव्यक्ती की जिची मुक्त चरे कोणतीही भरली तरी ती नेमकी एकच वस्तू दर्शवते. उदाहरणार्थ, ``$\emptyset$'' हे रिकामा संच दर्शवणारे पद आहे; ``$\{x\}$'' हे $x$ चा एकसदस्यीय संच दर्शवणारे पद आहे ($x$ काहीही असो); ``$x \cup y$'' हे $x$ आणि $y$ यांचा संयोग दर्शवणारे पद आहे (ते काहीही असोत).}

\begin{cor}
कोणत्याही पद $\tau(x)$ आणि कोणत्याही संच $A$ साठी पुढील संच अस्तित्वात असतो:
\[
	\Setabs{\tau(x)}{x \in A} = \Setabs{y}{(\exists x \in A)y = \tau(x)}.
\]
\end{cor}

\begin{proof}
$\tau$ हे \emph{पद} असल्याने $\forall x \lexists![y][\tau(x) = y]$. विशेषतः $(\forall x \in A)\lexists![y][\tau(x) = y]$. म्हणून प्रतिस्थापनावरून $\Setabs{y}{(\exists x \in A)\tau(x) = y}$ अस्तित्वात असतो.
\end{proof}
\noindent
यावरून या स्वयंसिद्धकाला ``प्रतिस्थापन'' हे नाव योग्य वाटते: संच $A$ दिला असता, $A$ च्या प्रत्येक सदस्याच्या जागी $\tau$ अंतर्गत त्याची प्रतिमा ठेवून $\Setabs{\tau(x)}{x \in A}$ हा नवा संच तयार करता येतो. फलनाखालील संचाच्या प्रतिमेचा संकेत अनुसरून $\Setabs{\tau(x)}{x \in A}$ साठी $\funimage{\tau}{A}$ असेही लिहिता येईल.

पण येथे महत्त्वाची बाब म्हणजे $\tau$ हे \emph{पद} आहे. \olref[sfr][fun][rel]{sec} मधील अर्थाने, म्हणजे क्रमित जोड्यांचा विशिष्ट संच या अर्थाने, ते \emph{फलनाचे} नाव असणे आवश्यक नाही. कारण $f$ हे त्या अर्थाने फलन असेल आणि $A\subseteq\dom{f}$,
तर $\funimage{f}{A}=\Setabs{f(x)}{x\in A}$ हा संच $\ran{f}$ चा
केवळ विशिष्ट उपसंच आहे. मनमाना $A$ घेतल्यास $f(x)$ हा संकेत
फक्त $A\cap\dom{f}$ मधील $x$ साठी वापरता येतो; प्रतिमा
क्रमित जोड्यांच्या आलेखावरून त्याच बंधनाने घेतली जाते. त्याचे अस्तित्व तर फक्त $\Zminus$ च्या स्वयंसिद्धकांवरून आधीच मिळते.\footnote{$\Setabs{y\in\bigcup\bigcup f}{(\exists x\in A)\tuple{x,y}\in f}$ पहा.} याउलट, प्रतिस्थापन हे आपल्या स्वयंसिद्धकांमध्ये \emph{सामर्थ्यवान} भर घालते; हे \olref[sth][replacement][]{chap} मध्ये दिसेल.

\end{document}
